\chapter{Машина из живых нейронов}
\chaptermark{Машина из живых нейронов}
\label{sec:livingmachine}

\section{Стенд с открытой схемой}

В главе~\ref{sec:debugging} мы отлаживали машину, у которой нет схемы: щуп слышит тысячу нейронов из восьмидесяти миллиардов, и всё остальное приходится угадывать. Инженер в таком положении сделал бы простую вещь --- собрал бы маленький кусок схемы на макетной плате, где видна каждая деталь. С нейронами это тоже можно сделать.

Нейроны коры разделяют и выращивают на чипе с решёткой электродов --- от нескольких десятков до десятков тысяч. Через несколько недель клетки прорастают отростками и соединяются в сеть из тысяч или сотен тысяч нейронов, в основном \nt{ГЛУТ} и примерно на пятую часть \nt{ГАМК}. Каждый электрод умеет и слушать, как щуп главы~\ref{sec:debugging}, и подавать ток, как тестовый сигнал. Вторая разновидность стенда --- \emph{органоиды}: трёхмерные комочки нервной ткани размером около миллиметра, выращенные из стволовых клеток человека~\cite{Lancaster2013}. На таких стендах живые сети работают месяцами; есть площадки, где опыты на органоидах можно ставить удалённо, по сети, больше ста дней подряд~\cite{Jordan2024}. Эту область называют \emph{биологическими вычислениями} или \emph{органоидным интеллектом}~\cite{Smirnova2023}.

У стенда два важных отличия от мозга. Он крошечный: в нём в сто тысяч раз меньше нейронов. И в нём нет широковещательных каналов: нейронов \nt{ДОФА}, \nt{СЕРО}, \nt{НОРА} и \nt{АЦЕТ} в культуре из коры нет. Это машина, у которой есть шина данных и управление потоком, но нет регистров управления. Посмотрим, на что она способна.

\section{Что сеть делает сама}

Оставленная в покое, культура не молчит и не шумит равномерно. Каждые несколько секунд почти вся сеть вспыхивает разом --- \emph{залпом} --- и снова затихает~\cite{Wagenaar2006}. Для инженера это знакомая картина: усилитель с сильной положительной обратной связью и без нагрузки превращается в генератор.

Механизм мы уже знаем. Сильные взаимные связи \nt{ГЛУТ}-нейронов запускают лавину, как в главе~\ref{sec:gaba} при нарушении условий утверждения~\ref{prop:ei}. Лавина быстро истощает запас медиатора в синапсах~\eqref{eq:depression}, и сеть замолкает. Запас восстанавливается с постоянной времени $\tau_{\text{вос}}$, и когда восстановится доля $x^*$, достаточная для новой лавины, следует новый залп. Интервал между залпами
\begin{empheq}[box=\keybox]{equation}
T_{\text{залп}} \;\approx\; \tau_{\text{вос}}\,\ln\frac{1}{1-x^*} .
\label{eq:burstinterval}
\end{empheq}
При $\tau_{\text{вос}}=0.8\,\text{с}$ из главы~\ref{sec:code} и $x^*=0.9$ получается около двух секунд, при $x^*=0.99$ --- около четырёх, то есть те самые несколько секунд. Это генератор релаксационного типа, родственник генератора ритма из главы~\ref{sec:muscles}: там его перезаряжала усталость групп, здесь --- запас медиатора.

Залпы --- то, что сеть делает, когда ей нечего делать. В живом мозге похожая картина бывает в глубоком сне (глава~\ref{sec:sleep}) и в развивающемся мозге, до того как в него пошли настоящие входы. Сходство это наводит на мысль, но одинаковость механизмов --- гипотеза.

\section{Как учить сеть без учителя-дофамина}

Раз \nt{ДОФА} в культуре нет, сигнал «лучше или хуже» приходится подавать самим --- электродами. Опыты показывают, что сеть на стенде действительно меняет ответы, и самое интересное в них --- \emph{чем} её учат.

\emph{Учитель, который замолкает.} Сеть стимулируют через один электрод раз в секунду-три, пока на другом электроде не появится нужный ответ через $50\pm10\ms$ после стимула. Как только он появился, стимуляцию выключают на пять минут. После нескольких таких циклов нужный ответ появляется на стимул сразу~\cite{Shahaf2001}. Награда здесь --- тишина: стимул расшатывает сеть, а прекращение стимула оставляет её в том состоянии, которое дало правильный ответ. Это спуск по градиенту через случайные отклонения из главы~\ref{sec:dofa} (утверждение~\ref{prop:gradient}), где роль ошибки играет сам факт тряски: трясут, пока не станет правильно.

\emph{Сеть, которая играет в теннис.} В опыте на стенде с плотной решёткой электродов около миллиона нейронов подключили к компьютерной игре в теннис с одной ракеткой~\cite{Kagan2022}. Положение мяча подавалось током на «чувствительные» электроды: где мяч --- кодом местом, как далеко --- частотой. Активность в двух «двигательных» областях двигала ракетку вверх или вниз. Правило обучения было необычным. Если ракетка отбила мяч, сеть получала короткий \emph{предсказуемый} стимул; если промахнулась --- несколько секунд \emph{непредсказуемого} случайного тока. За двадцать минут игры серии отбитых мячей стали длиннее, чем у контрольных культур без обратной связи. Подробный разбор этого опыта --- в главе~\ref{sec:dishbrain}.

Авторы объяснили результат гипотезой, что сеть действует так, чтобы её вход был предсказуемым~\cite{Friston2010}. Это та же идея, что экономия импульсов в главе~\ref{sec:code} и предсказание кадра в главе~\ref{sec:recognition}: машина тратит меньше, когда вход ожидаем. Но и сам опыт, и особенно слова авторов о «разумности» культуры, вызвали резкую критику: эффект невелик, а объяснить его можно и проще~\cite{Balci2023}. Честная оценка такая: культура на стенде меняет поведение под действием обратной связи --- это измерено, а что она учится именно предсказывать свой вход --- гипотеза.

\emph{Сеть, которая водит тело.} Похожим образом культуру подключали к простому виртуальному «телу» и учили его двигаться к цели, подбирая пространственно-временной рисунок тренирующих стимулов~\cite{Bakkum2008}. Ни в одном из этих опытов нет дофамина: учителем служит рисунок тока, который выбирает инженер. Что меняется, когда учителем становится программа или сам дофамин, --- в разделах~\ref{sec:dishdopamine}--\ref{sec:cartpole}.

\begin{figure}[t]
\centering
\begin{tikzpicture}[>=Stealth,
  blk/.style={draw,very thick,rounded corners=2pt,align=center,font=\footnotesize,minimum height=0.8cm,fill=blue!5}]
% (a) closed loop
\node[font=\small] at (1.5,4.3) {(а) замкнутая петля};
\node[blk,text width=2.6cm] (G) at (1.5,3.3) {игра: мяч и ракетка};
\draw[very thick,fill=orange!10,rounded corners=3pt] (0.5,0.2) rectangle (2.5,1.6);
\foreach \x in {0.7,1.0,...,2.35}{\foreach \y in {0.4,0.7,1.0,1.3}{\fill[gray!60] (\x,\y) circle (0.04);}}
\draw[->,thick,blue!60!black] (G.west) -| (-0.5,0.9) -- (0.5,0.9);
\node[font=\scriptsize,blue!60!black,anchor=east] at (-0.6,2.1) {мяч $\to$ ток};
\draw[->,thick,red!70!black] (2.5,0.9) -- (3.5,0.9) |- (G.east);
\node[font=\scriptsize,red!70!black,anchor=west] at (3.6,2.1) {импульсы $\to$ ракетка};
\node[font=\scriptsize] at (1.5,-0.2) {нейроны на решётке электродов};
\node[font=\scriptsize,align=center] at (1.5,-0.85) {отбил: предсказуемый стимул\\промах: случайный ток};
% (b) reservoir
\begin{scope}[xshift=7.9cm]
\node[font=\small] at (1.6,4.3) {(б) резервуар};
\node[blk,text width=1.1cm] (X) at (-0.4,1.0) {вход\\$u(t)$};
\draw[very thick,fill=orange!10] (1.6,1.0) circle (0.8);
\foreach \a/\r in {20/0.35,80/0.5,150/0.3,210/0.55,280/0.4,330/0.2,110/0.15}{\fill[red!60!black] ({1.6+\r*cos(\a)},{1.0+\r*sin(\a)}) circle (0.05);}
\node[font=\scriptsize] at (1.6,-0.2) {органоид: не обучается};
\node[blk,text width=1.8cm,fill=green!8] (W) at (4.0,1.0) {считывание\\$W_{\text{вых}}$};
\draw[->,thick] (X) -- (0.8,1.0);
\draw[->,thick] (2.4,1.0) -- node[above,font=\scriptsize] {$\mathbf r(t)$} (W);
\draw[->,thick] (W) -- (4.0,2.3) node[above,font=\scriptsize] {$y(t)$};
\node[font=\scriptsize,green!40!black] at (4.0,0.3) {обучается};
\end{scope}
\end{tikzpicture}
\caption{Два способа заставить живую сеть вычислять. (а) Замкнутая петля: положение мяча подаётся током, импульсы двигательных областей двигают ракетку, а предсказуемый или случайный стимул после удара служит учителем. (б) Резервуар~\eqref{eq:reservoir}: органоид не обучают вовсе, он только расщепляет вход на множество нелинейных откликов с памятью; учится одно линейное считывание снаружи.}
\label{fig:livingmachine}
\end{figure}

\section{Живой резервуар}

Есть и другой способ использовать живую сеть: не учить её вовсе. В технике он называется \emph{вычислением на резервуаре}~\cite{Maass2002,JaegerHaas2004}. Берут готовую нелинейную систему с памятью --- любую, лишь бы её отклик на вход был богатым и зависел от недавнего прошлого. Вход $u(t)$ раскачивает её, получается вектор откликов $\mathbf r(t)$, и обучается только линейное считывание
\begin{empheq}[box=\keybox]{equation}
y(t) \;=\; W_{\text{вых}}\,\mathbf r(t) ,
\label{eq:reservoir}
\end{empheq}
веса которого подбираются правилом наименьших квадратов~\eqref{eq:lms} из главы~\ref{sec:motorlearning}. Резервуар делает то, что делали маленькие \nt{ГЛУТ}-нейроны главы~\ref{sec:motorlearning}: расщепляет вход на длинный вектор, в котором нужные классы становятся разделимы плоскостью (глава~\ref{sec:dendrites}).

Органоид на решётке электродов подошёл на роль такого резервуара. Звуки речи, поданные на него рисунками тока, после обучения одного линейного считывания различались по говорящему с точностью около $78\%$~\cite{Cai2023}. Результат честный и полезный, но важно понимать, где в нём «ум»: органоид даёт нелинейность и затухающую память, а всё обучение сделано снаружи, в кремнии (рис.~\ref{fig:livingmachine}).

\section{Что стенд говорит о машине}

\emph{Энергия.} По оценке~\eqref{eq:energy} импульс стоит около $10^{-10}$~Дж. Культура из миллиона нейронов при частоте в импульс в секунду тратит на передачу сигналов
\begin{empheq}[box=\keybox]{equation}
P \;\approx\; 10^{6}\times1\,\text{с}^{-1}\times10^{-10}\,\text{Дж} \;=\; 10^{-4}\,\text{Вт},
\label{eq:dishpower}
\end{empheq}
десятую долю милливатта. Электроника, которая стимулирует, записывает и обрабатывает эти импульсы, тратит на порядки больше. Как и в главе~\ref{sec:debugging}, узкое место --- не вычисление, а связь с ним.

\emph{Недостающие детали.} Стенд показывает, сколько умеет шина \nt{ГЛУТ} с управлением потоком \nt{ГАМК} без регистров управления: самоорганизоваться, генерировать залпы, менять ответы под действием обратной связи и служить хорошим резервуаром. Чего она без них не умеет, --- сама решить, что считать успехом. Этот сигнал на стенде подаёт инженер, а в мозге, как мы видели в главах~\ref{sec:dofa}--\ref{sec:acet}, --- широковещательные каналы.

\emph{Этика.} Органоиды выращивают из клеток человека, и чем сложнее они становятся, тем серьёзнее вопрос, может ли такая ткань что-то чувствовать. Инженерный взгляд подсказывает частичный ответ: боль в главе~\ref{sec:pain} --- это не сигнал одного датчика, а целая система тревоги с воротами, регулятором громкости и широковещательными каналами, которых на стенде нет. Но это рассуждение, а не доказательство, и сама область предлагает вести этическую оценку вместе с опытами, с самого начала~\cite{Smirnova2023}.

\begin{keyblock}
\begin{proposition}[Живая сеть на стенде]
\label{prop:livingmachine}
Сеть из \nt{ГЛУТ}- и \nt{ГАМК}-нейронов на решётке электродов сама порождает залпы~\eqref{eq:burstinterval}, меняет ответы под действием обратной связи и может служить резервуаром~\eqref{eq:reservoir} для обученного снаружи считывания. Всё это измерено. Что такая сеть учится по какому-то общему правилу, например предсказывать свой вход, --- гипотеза, а громкие слова о её «разумности» большинство специалистов не принимает.
\end{proposition}
\end{keyblock}

\section{Дофамин по вспышке света}
\label{sec:dishdopamine}

Единственный известный нам прямой опыт, где культуре дали дофамин как учителя, поставлен на платформе из органоидов, к которой исследователи подключаются удалённо~\cite{Jordan2024}. В жидкость, омывающую органоиды, добавили дофамин в светочувствительной «клетке»: связанная молекула не действует на рецепторы. Концентрация клеточного дофамина --- $1$~мМ. Когда нужно наградить сеть, на неё светят ультрафиолетом: клетка распадается, и дофамин выходит в объём.

Петля устроена так. Один и тот же стимул подаётся раз за разом; если он вызвал больше импульсов, чем прежде, включается вспышка и высвобождается дофамин. За $240$ повторов число вызванных импульсов в целом росло. Авторы сами пишут, что по одному такому опыту нельзя заключить, что рост вызван дофамином~\cite{Jordan2024}.

Для инженера это первая попытка собрать в чашке правило трёх множителей~\eqref{eq:threefactor}: активность отправителя и получателя оставляет след, а общий сигнал решает, закрепить ли изменение. Но сигнал здесь бывает только положительным: награда есть или её нет, отрицательной ошибки $\delta<0$ установка не выдаёт. Дофамин к тому же расплывается и откачивается медленно, как концентрация в~\eqref{eq:lowpass}, так что частые вспышки могут просто поднять его общий уровень. Чтобы отличить обучение от этого, нужны два контроля: столько же вспышек в случайные моменты, не связанные с ответом сети, и те же вспышки без клеточного дофамина. И нужна проверка после отдыха: осталось ли изменение, когда дофамин ушёл.

\section{Дофамин учит кремний}
\label{sec:biohybrid}

В обратном опыте живые клетки учат электронику~\cite{Keene2020}. Клетки, выделяющие дофамин, выращивают в микроканале над органическим нейроморфным элементом --- транзистором, проводимость канала которого играет роль веса синапса. Дофамин, вышедший из клеток, окисляется у электрода элемента, и это меняет проводимость. Устройство повторяет и механизм уборки медиатора из щели, поэтому вес можно не только изменить надолго, но и вернуть обратно.

Схемотехнически это интегратор с утечкой по химическому входу: вес копит поток дофамина, а «откачка» задаёт, как быстро он забывается, --- та же RC-цепь, что и~\eqref{eq:lowpass}. Главное здесь направление связи. Щуп из главы~\ref{sec:debugging} не видит широковещательных регистров, потому что они химические. Этот элемент читает именно химический регистр и превращает его в электрический вес. Для интерфейсов мозг---компьютер это путь к датчику, который слышит не только шину данных, но и регистры управления.

\section{Органоиды со своими дофаминовыми нейронами}
\label{sec:midbrainorg}

Из стволовых клеток человека умеют выращивать органоиды, похожие на ту область мозга, где сидят \nt{DOPA}-нейроны. В них есть работающие дофаминовые нейроны, которые выделяют дофамин~\cite{Jo2016}. Такие органоиды выращивают в основном для изучения болезней. Опытов, где их дофамин служил бы учителем для другой сети, мы не нашли.

Для машины из живых нейронов это недостающая деталь. Стенд из главы~\ref{sec:livingmachine} --- шина данных и управление потоком без регистров управления; здесь источник широковещательного сигнала уже есть. Не хватает критика: сети, которая вычисляла бы ошибку предсказания~\eqref{eq:ctd} и управляла бы этими нейронами. Соединить органоид коры с таким органоидом так, чтобы дофамин выделялся по ошибке, --- открытая задача, и пока это гипотеза, а не опыт.

\section{Органоид балансирует шестом без дофамина}
\label{sec:cartpole}

Самый полный из недавних опытов обходится без дофамина совсем~\cite{Robbins2026}. Органоиды коры мыши подключили в замкнутую петлю к задаче «шест на тележке»: активность органоида двигает тележку, положение шеста подаётся обратно стимуляцией. Между попытками органоид получает короткие высокочастотные учебные стимулы. Какие именно --- выбирает программа обучения с подкреплением.

Результаты измерены:
\begin{itemize}
\item у большинства органоидов учебные сигналы, выбранные программой, дали лучший результат, чем случайно выбранные сигналы или отсутствие сигнала;
\item после $45$ минут отдыха улучшение не сохранялось;
\item блокада рецепторов глутамата обоих типов, AMPA и NMDA, отменяла улучшение.
\end{itemize}

Последний пункт говорит, где живёт выученное: в \nt{GLUT}-синапсах, и через тот рецептор, который в разделе~\ref{sec:weightstore} работает как детектор совпадения входа и выхода. Второй пункт показывает, чего не хватает: быстрое изменение есть, закрепления нет, как у быстрого ученика из главы~\ref{sec:motorlearning} без медленного. А роль учителя изменилась: инженера, который выбирает рисунок тока, заменила программа, но учитель по-прежнему стоит снаружи чашки.

Среди более ранних опытов обучения культур на электродных матрицах мы также не нашли ни одного, где использовался бы дофамин: учебным сигналом везде служила электрическая стимуляция.

\section{Кто учит культуру}

\begin{table}[t]
\centering
\footnotesize
\setlength{\tabcolsep}{4pt}
\renewcommand{\arraystretch}{1.15}
\begin{tabular}{>{\raggedright}p{2.7cm}>{\raggedright}p{2.5cm}>{\raggedright}p{3.2cm}>{\raggedright\arraybackslash}p{4.4cm}}
\toprule
Опыт & Кто учит & Учебный сигнал & Что известно \\
\midrule
остановка стимуляции при нужном ответе & инженер & конец стимуляции & ответ закрепляется --- измерено \\
DishBrain (гл.~\ref{sec:dishbrain}) & инженер & предсказуемый или случайный ток & значимый рост розыгрышей за сеанс только при полной обратной связи, при тишине эффект меньше --- измерено; от сеанса к сеансу неустойчив; причина обсуждается \\
шест на тележке & программа обучения с подкреплением & выбранные программой высокочастотные стимулы & лучше случайных, но не сохраняется после отдыха --- измерено \\
дофамин по вспышке & правило «больше импульсов --- награда» & дофамин, освобождённый светом & рост импульсов --- наблюдение; роль дофамина не доказана \\
биогибридный синапс & живые клетки & дофамин от клеток & долгое изменение и возврат веса элемента --- измерено \\
органоиды с \nt{DOPA}-нейронами & --- & --- & источник дофамина есть; как учитель не использован \\
\bottomrule
\end{tabular}
\caption{Кто и чем учит живые нейроны на чипе.}
\label{tab:dishteachers}
\end{table}

Во всех опытах, где учится сама культура, учитель пока снаружи (таблица~\ref{tab:dishteachers}): инженер, программа или правило, заданное инженером. Дофамин вошёл в чашку лишь однажды, и его роль ещё предстоит доказать. Собрать в чашке настоящий широковещательный канал --- с критиком, ошибкой и следом, как в главе~\ref{sec:dofa}, --- следующий шаг, которого никто ещё не сделал.


\section{Итог}

\begin{table}[t]
\centering
\footnotesize
\setlength{\tabcolsep}{4pt}
\renewcommand{\arraystretch}{1.15}
\begin{tabular}{>{\raggedright}p{2.8cm}>{\raggedright}p{4.2cm}>{\raggedright\arraybackslash}p{5.8cm}}
\toprule
Опыт & Что делает живая сеть & Что известно \\
\midrule
сеть без входа & залпы каждые несколько секунд~\eqref{eq:burstinterval} & измерено; механизм через истощение синапсов --- общепринятая модель \\
учитель, который замолкает & нужный ответ закрепляется, когда стимуляцию выключают & измерено \\
игра в теннис & серии отбитых мячей удлиняются при предсказуемой обратной связи & изменение поведения измерено; величина эффекта и объяснение предсказанием входа спорны \\
виртуальное тело & движение к цели при подобранном рисунке стимулов & измерено \\
резервуар & нелинейность и память для считывания~\eqref{eq:reservoir} & измерено; обучение --- снаружи, в кремнии \\
энергия & около $10^{-4}$~Вт на миллион нейронов~\eqref{eq:dishpower} & оценка по~\eqref{eq:energy} \\
\bottomrule
\end{tabular}
\caption{Что показали машины из живых нейронов.}
\label{tab:livingmachine}
\end{table}

Машина из живых нейронов (таблица~\ref{tab:livingmachine}) --- это стенд, на котором видно, что умеют шина данных и управление потоком без регистров управления. Они умеют многое, но не умеют сами решить, что хорошо. На стенде эту роль играет инженер со своим рисунком тока; в мозге --- те немногие широковещательные провода, которым посвящена середина книги.

\section{Задачи}

\begin{exercise}[\lvlA]
\label{ex:21:1}
\emph{Интервал между залпами.} (а)~Выведите~\eqref{eq:burstinterval}, считая, что залп истощает запас до нуля, а затем запас $x$ восстанавливается по закону $\tau_{\text{вос}}\,\dd x/\dd t=1-x$. Обобщите формулу на случай, когда залп оставляет долю $x_0$.
(б)~Найдите $T_{\text{залп}}$ при $\tau_{\text{вос}}=0.8$~с и $x^*=0.9$, $0.99$ и чувствительность $\dd T_{\text{залп}}/\dd x^*$ в обоих случаях.
(в)~Пусть $x^*$ от залпа к залпу случайно меняется на $\pm0.005$ около $0.99$. В каком диапазоне лежат интервалы? Что это говорит о регулярности залпов?
\end{exercise}

\begin{exercise}[\lvlA]
\label{ex:21:2}
\emph{Энергия стенда.} (а)~Проверьте~\eqref{eq:dishpower} и, применив ту же оценку ко всему мозгу ($8.6\cdot10^{10}$ нейронов), сравните с~\eqref{eq:energy}.
(б)~Стенд записывает $1024$ электрода с частотой $20$~кГц по $10$ бит на отсчёт. Каков поток данных? Сколько тратит его передача при $1$~нДж на бит и во сколько раз это больше мощности самой культуры?
\end{exercise}

\begin{exercise}[\lvlB]
\label{ex:21:3}
\emph{Проектирование: подавить залпы.} Залпы мешают культуре работать резервуаром: каждый стирает её память. Предлагается подавать через много электродов редкую фоновую стимуляцию, так что каждый синапс выбрасывает медиатор с частотой $r_b$. Используйте модель истощения главы~\ref{sec:code}: $\tau_{\text{вос}}\,\dd x/\dd t=1-x-U r\tau_{\text{вос}}\,x$ с $U=0.5$, $\tau_{\text{вос}}=0.8$~с.
(а)~Найдите установившийся запас $x_{\text{уст}}(r_b)$ и наименьшую $r_b$, при которой новая лавина невозможна, для $x^*=0.9$ и $0.99$.
(б)~Какова цена этого решения для передачи сигнала через синапс?
(в)~Почему стимуляция должна быть распределена по многим электродам, а не подана через один?
\end{exercise}

\begin{exercise}[\lvlB]
\label{ex:21:4}
\emph{Память резервуара не больше числа нейронов.} Резервуар с $N$ выходами $\mathbf r(t)$ возбуждается стационарным случайным входом $u(t)$ с нулевым средним, единичной дисперсией и независимыми значениями в разные моменты. Для каждой задержки $k\ge0$ линейное считывание~\eqref{eq:reservoir} $y_k(t)=\mathbf w_k\cdot\mathbf r(t)$ подбирается так, чтобы максимизировать квадрат коэффициента корреляции $m_k$ между $y_k(t)$ и $u(t-k)$. Докажите, что ёмкость памяти
\begin{equation*}
\mathrm{MC}=\sum_{k=0}^{\infty}m_k\le N .
\end{equation*}
Указание: рассмотрите случайные величины как векторы со скалярным произведением $\langle X,Y\rangle=\mathbb E[XY]$; что можно сказать о системе $\{u(t-k)\}_k$?
\end{exercise}

\begin{exercise}[\lvlB]
\label{ex:21:5}
\emph{Моделирование: резервуар в кремнии.} Напишите на чистом Python резервуар из $N=30$ элементов: $\mathbf r(t+1)=\tanh\bigl(W\mathbf r(t)+\mathbf v\,u(t)+\mathbf c\bigr)$; $W$ --- случайная гауссова матрица, отнормированная на спектральный радиус $0.9$ (степенной метод); $v_i\sim U(-0.5,0.5)$, $c_i\sim U(-0.2,0.2)$. Считывание --- гребневая регрессия (метод наименьших квадратов с регуляризатором $10^{-6}$ на отсчёт) по $1500$ шагам после $100$ шагов разгона, проверка на следующих $1000$.
(а)~При $u(t)\sim U(-1,1)$ найдите $m_k$ для $k=0,\dots,89$ и $\mathrm{MC}$. Повторите без $\tanh$ (линейный резервуар) при спектральном радиусе $0.9$ и $0.99$. Проверьте неравенство задачи~\ref{ex:21:4}.
(б)~При $u(t)=\pm1$ обучите считывание выдавать $y(t)=u(t-1)\,u(t-2)$ (исключающее ИЛИ). Сравните с линейным считыванием прямо с отводов $u(t)$, $u(t-1)$, $u(t-2)$ и с резервуаром без смещений ($\mathbf c=0$). Объясните последний результат.
\end{exercise}

\begin{exercise}[\lvlB]
\label{ex:21:6}
\emph{Проектирование считывания для живого резервуара.} С органоида записывают $1024$ канала, признаки --- счёт импульсов канала за окно. Для обучения есть $P=240$ записей с двумя классами.
(а)~Пользуясь формулой задачи~\ref{ex:19:3}, объясните, почему считывание по всем $1024$ признакам даст стопроцентную точность на обучающих данных при любых, даже случайных, метках.
(б)~Найдите наибольшее число признаков $n$ (после их объединения в группы), при котором случайная разметка $240$ записей разделима с вероятностью не больше $1\%$.
(в)~На $100$ проверочных записях получена точность $78\%$. Насколько она надёжно отличается от случайного угадывания?
\end{exercise}

\begin{exercise}[\lvlC]
\label{ex:21:7}
\emph{Исследование: учитель без широковещательного канала.} По утверждению~\ref{prop:livingmachine} стенд не умеет сам решить, что считать успехом; эту роль играет рисунок тока.
(а)~Предложите протокол стимуляции, который реализует на стенде трёхфакторное правило главы~\ref{sec:dofa}: что играет роль следа, что --- роль ошибки $\delta$, и как подать её всем синапсам сразу через $8$--$1024$ электродов?
(б)~Сравните ожидаемую скорость обучения с «учителем, который замолкает» (опыт с тишиной как наградой), используя~\eqref{eq:perturbation} и оценку шума считывания~\eqref{eq:rateerror}.
(в)~Как отличить на стенде обучение весов в сети от сдвига её текущего состояния? Предложите контроль, в котором считывание заморожено, и назовите результат, который опроверг бы гипотезу об обучении.
\end{exercise}
